\section{A high-concentration Hankel cutoff}
\label{sec:hankel-high-concentration}

The following selection is different from the fixed threshold
$1/2$ used in Section~\ref{sec:hankel-selection}. This difference is
part of the definition: the present cutoff retains only modes whose
concentration error is at most a prescribed negative power of
$\Lambda$. All constants in this section may depend on the fixed
Hankel parameter $c>0$.

On $L^2((0,\infty),dx)$ let
\[
 P_\Lambda=1_{(0,\Lambda)},\qquad
 B_\Lambda=\mathcal H_cP_\Lambda\mathcal H_c,\qquad
 A_\Lambda=P_\Lambda B_\Lambda P_\Lambda,
 \qquad F_\Lambda=P_\Lambda\mathcal H_cP_\Lambda.
\]
Thus $A_\Lambda=F_\Lambda^2$. Fix $M>0$, put
$\delta_\Lambda=\Lambda^{-M}$, and, for sufficiently large $\Lambda$,
define
\begin{equation}
 \label{eq:hankel-high-cutoff}
 \Pi_{\Lambda,M}=1_{[1-\delta_\Lambda,1]}(A_\Lambda),
 \qquad
 W_{\Lambda,M}=B_\Lambda
       (A_\Lambda|_{\operatorname{Ran}\Pi_{\Lambda,M}})^{-1}.
\end{equation}
The inverse in the last expression acts only on the selected space.
As before, $H_D^s$ denotes the Sobolev norm after the unitary change
of variables $Uf(t)=e^{t/2}f(e^t)$.

\begin{proposition}[Size, extension, and graph tail]
\label{prop:hankel-high-cutoff}
The selection in \eqref{eq:hankel-high-cutoff} is finite and
\begin{equation}
 \label{eq:hankel-high-rank}
 \operatorname{rank}\Pi_{\Lambda,M}
                 \sim\frac{2c\Lambda}{\pi}.
\end{equation}
Its selected eigenvalues are at least $1-\Lambda^{-M}>0$;
every omitted positive eigenvalue is strictly less than
$1-\Lambda^{-M}<1$. For $\psi$ in its range,
\begin{align}
 P_\Lambda W_{\Lambda,M}\psi&=\psi,\label{eq:high-right-inverse}\\
 \|W_{\Lambda,M}\psi\|_2^2
        &\leq(1-\delta_\Lambda)^{-1}\|\psi\|_2^2,\label{eq:high-extension-norm}\\
 \|(1-P_\Lambda)W_{\Lambda,M}\psi\|_2^2
        &\leq\frac{\delta_\Lambda}{1-\delta_\Lambda}
                                      \|\psi\|_2^2.\label{eq:high-raw-tail}
\end{align}
For every fixed $0\leq\sigma<s<1/2$,
\begin{equation}
 \label{eq:high-graph-tail}
 \|(1-P_\Lambda)W_{\Lambda,M}\psi\|_{H_D^\sigma}
 \leq C_{c,M,s,\sigma}
       \Lambda^{\sigma-\frac M2(1-\sigma/s)}\|\psi\|_2.
\end{equation}
In particular, any prescribed negative power on the right is
obtained by choosing a sufficiently large fixed $M$.
\end{proposition}

\begin{proof}
Finiteness follows from $\Tr A_\Lambda<\infty$. Since
$\Pi_{\Lambda,M}\leq1_{[1/2,1]}(A_\Lambda)$ for large $\Lambda$,
all the raw estimates in
Proposition~\ref{prop:hankel-raw-estimates} apply uniformly to this
smaller space. Directly,
\[
 \|W_{\Lambda,M}\psi\|_2^2
     =\langle\psi,A_\Lambda^{-1}\psi\rangle,
 \qquad P_\Lambda W_{\Lambda,M}\psi=\psi.
\]
Subtract $\|\psi\|_2^2$ from the first identity to obtain
\eqref{eq:high-raw-tail}.

For the global Sobolev bound, the finite Hankel transform commutes
with its square and therefore preserves every spectral subspace
used here. It is invertible on the selected space and
\[
 \|F_\Lambda^{-1}\Pi_{\Lambda,M}\|
                       \leq(1-\delta_\Lambda)^{-1/2}.
\]
Moreover
\begin{equation}
 \label{eq:high-extension-hankel}
 W_{\Lambda,M}\psi
            =\mathcal H_cP_\Lambda F_\Lambda^{-1}\psi.
\end{equation}
Indeed, $F_\Lambda A_\Lambda^{-1}=F_\Lambda^{-1}$ on this
subspace. The identity
$\mathcal H_cD_r=D_{1/r}\mathcal H_c$ shows that $\mathcal H_c$
is an isometry on every $H_D^s$. Apply the earlier raw estimate
to $F_\Lambda^{-1}\psi$ to get
\[
 \|W_{\Lambda,M}\psi\|_{H_D^s}
                         \leq C_{c,s}\Lambda^s\|\psi\|_2.
\]
Multiplication by $1_{(\log\Lambda,\infty)}$ is bounded on
$H^s(\mathbb R)$, uniformly in the endpoint, for $0<s<1/2$.
One direct proof uses the fractional difference quotient formula:
the cross-boundary contribution is a constant multiple of
$\int_0^\infty |f(L+u)|^2u^{-2s}du$, which is bounded by the
$H^s$ norm by the one-dimensional fractional Hardy inequality.
For completeness, apply the difference quotient formula on
$I_j=(2^{-j-1},2^{-j})$, $j\geq0$, and let $m_j$ be the mean
of $f(L+\cdot)$ on $I_j$. The oscillation terms satisfy
\[
 \sum_{j\geq0}|I_j|^{-2s}
       \int_{I_j}|f(L+u)-m_j|^2du
                    \leq C_s[f]_{\dot H^s}^2.
\]
The difference of two adjacent means is bounded by their
cross-interval difference quotient integral times
$C_s|I_j|^{2s-1}$. Put $\rho=2^{-(1-2s)}<1$.
The elementary discrete inequality
\[
 \sum_{j\geq0}\rho^j|m_j|^2
 \leq C_\rho\left(|m_0|^2+
                \sum_{j\geq0}\rho^j|m_{j+1}-m_j|^2\right)
\]
follows by writing each $m_j$ as a partial sum and applying
the $\ell^1*\ell^2$ convolution bound to the sequence
$\rho^{j/2}$. This controls the mean terms by the same
fractional seminorm and an $L^2$ term. The region $u\geq1$
is controlled directly by $\|f\|_2^2$. These estimates prove
the asserted Hardy inequality, and explain the restriction
$2s<1$.
Consequently the exterior part also has $H_D^s$ norm
$O_{c,s}(\Lambda^s\|\psi\|_2)$. Interpolate this estimate with
\eqref{eq:high-raw-tail}, using the Fourier definition of the
Sobolev norms and H\"older's inequality. The resulting powers are
$\Lambda^{-M(1-\sigma/s)/2}\Lambda^\sigma$, as asserted.

It remains to prove the lower rank asymptotic; a transition-width
theorem is unnecessary. Fix $0<a<1$, and choose a smooth function
$h$, equal to one on $[0,a]$, supported in $[0,b]$ for some $b<1$,
with $0\leq h\leq1$. Put $p=h^2$ and
\[
 K^p_\Lambda=P^p_\Lambda B^p_\Lambda P^p_\Lambda,
 \quad P^p_\Lambda=p(x/\Lambda),
 \quad B^p_\Lambda=\mathcal H_cp(x/\Lambda)\mathcal H_c.
\]
Then $0\leq K^p_\Lambda\leq1$. The strict-margin Fourier argument
in Proposition~\ref{prop:hankel-interior-approximation} gives
\[
 \|(1-B_\Lambda)P^p_\Lambda(B^p_\Lambda)^{1/2}\|_{2\to2}
                                         =O_N(\Lambda^{-N})
\]
for every $N$, since $(B^p_\Lambda)^{1/2}
=\mathcal H_ch(x/\Lambda)\mathcal H_c$ has the same strict
frequency margin. In addition,
\[
 \|(B^p_\Lambda)^{1/2}P^p_\Lambda\|_{\rm HS}^2
    =\Tr(P^p_\Lambda B^p_\Lambda P^p_\Lambda)
    \leq\Tr A_\Lambda=O_c(\Lambda).
\]
The inequality follows from the diagonal kernel, since both
soft factors lie between zero and their corresponding hard
cutoffs. Hence
$\|(1-B_\Lambda)K^p_\Lambda\|_{\rm HS}=O_N(\Lambda^{-N})$
after relabelling $N$.

For every vector $f$ supported in $(0,\Lambda)$ the omitted
spectral gap gives
\begin{equation}
 \label{eq:high-projection-gap}
 \|(1-\Pi_{\Lambda,M})f\|_2^2
 \leq\delta_\Lambda^{-1}\langle f,(1-A_\Lambda)f\rangle
 =\delta_\Lambda^{-1}\|(1-B_\Lambda)f\|_2^2.
\end{equation}
Apply this to the columns of $K^p_\Lambda$ and sum. Thus
\begin{equation}
 \label{eq:high-soft-HS}
 \|(1-\Pi_{\Lambda,M})K^p_\Lambda\|_{\rm HS}
                                      =O_N(\Lambda^{-N})
\end{equation}
for every fixed $N$. Since $K^p_\Lambda$ is a contraction,
\[
 \operatorname{rank}\Pi_{\Lambda,M}
 \geq\|\Pi_{\Lambda,M}K^p_\Lambda\|_{\rm HS}^2
 =\|K^p_\Lambda\|_{\rm HS}^2-o(1).
\]
Compress $K^p_\Lambda$ to $(0,a\Lambda)$. The time multiplier is
then the identity and its band multiplier is at least
$B_{a\Lambda}$, so this compression is at least $A_{a\Lambda}$.
For positive Hilbert--Schmidt operators, order implies order of
their squared traces: if $T=S+D$ with $S,D\geq0$, then
$\Tr T^2-\Tr S^2=2\Tr(SD)+\Tr D^2\geq0$.
It follows that
\[
 \operatorname{rank}\Pi_{\Lambda,M}
 \geq\Tr(A_{a\Lambda}^2)-o(1)
 =\frac{2ca\Lambda}{\pi}+O_{c,a}(\log\Lambda).
\]
Here the last equality uses the trace and variance estimates in
Proposition~\ref{prop:hankel-selection-size}.
On the other hand,
\[
 \operatorname{rank}\Pi_{\Lambda,M}
       \leq\frac{\Tr A_\Lambda}{1-\delta_\Lambda}
       =\frac{2c\Lambda}{\pi}+o(\Lambda).
\]
Let $a$ increase to one after taking the lower limit. This proves
\eqref{eq:hankel-high-rank}.
\end{proof}

\begin{proposition}[Interior approximation for the new cutoff]
\label{prop:hankel-high-interior}
For fixed $M>0$, the strict-margin and fixed-profile conclusions
of Proposition~\ref{prop:hankel-interior-approximation} remain
valid with $\Pi_{\Lambda,M}$ in place of $\Pi_\Lambda$:
for every $\sigma<1/2$ and every $N>0$,
\[
 \|(1-\Pi_{\Lambda,M})P^p_\Lambda B^p_\Lambda P^p_\Lambda
                          \|_{L^2\to H_D^\sigma}
                                  =O_{c,p,M,\sigma,N}(\Lambda^{-N}).
\]
The projected two-profile repair of the principal mean and value
constraints is also valid. Imposing these constraints changes
the rank by at most two, so it preserves
\eqref{eq:hankel-high-rank}.
\end{proposition}
\begin{proof}
The only change in the raw approximation proof is that the
omitted spectral gap $1/2$ is replaced by $\delta_\Lambda$.
Equation~\eqref{eq:high-projection-gap} therefore costs
$\Lambda^{M/2}$ in norm. Every strict-margin and fixed-profile
raw error has an arbitrarily large polynomial decay rate, so
this fixed loss is absorbed by choosing that rate larger.
The coarse $H_D^s$ bound remains $O(\Lambda^s)$, since the new
selected space is contained in the old one. The same interpolation
argument proves the stated graph estimate. The mean and value
of the projected repair profiles converge to their original
normalized values with every polynomial rate, by the same
Cauchy--Schwarz and ordinary trace estimates as before. Their
constraint matrix is consequently invertible for large $\Lambda$.
\end{proof}

\begin{corollary}[An abstract theta-tail criterion]
\label{cor:hankel-high-theta-tail}
Suppose a fixed theta multiplier $E$ obeys
\[
 \|Ef\|_2\leq C\|f\|_{H_D^\sigma},\qquad 0\leq\sigma<1/2.
\]
Let $\Pi^0_{\Lambda,M}$ be either the full high-concentration
selection or its subspace satisfying the two principal
constraints. Suppose that a lower-tail operator $T_{\Lambda,g}$
has the exact identity
\[
 T_{\Lambda,g}\psi
 =-V(g)1_{(-\infty,-\log\Lambda)}
                   E(1-P_\Lambda)W_{\Lambda,M}\psi
       \qquad(\psi\in\operatorname{Ran}\Pi^0_{\Lambda,M}),
\]
where $V(g)$ is bounded on the output $L^2$ space, uniformly in
$\Lambda$. If the theta Gram operator on the selected space
satisfies $\|G_\Lambda^{-1}\|=O(\Lambda^a)$ for some fixed
$a\geq0$, then
\begin{equation}
 \label{eq:high-normalized-tail}
 \|T_{\Lambda,g}\Pi^0_{\Lambda,M}G_\Lambda^{-1/2}\|_{\rm HS}^2
 \ll\Lambda^{a+1+2\sigma-M(1-\sigma/s)},
                 \qquad \sigma<s<1/2.
\end{equation}
In particular the normalized tail tends to zero whenever
$M(1-\sigma/s)>a+1+2\sigma$. With the crude exponent $a=2$,
it is enough to take $M>(3+2\sigma)/(1-\sigma/s)$.
For example, the choices
\[
 \sigma=3/8,\qquad s=7/16,\qquad M=28,
\]
give an $O(\Lambda^{-1/4})$ bound in
\eqref{eq:high-normalized-tail} when $a=2$.
\end{corollary}
\begin{proof}
The assumed identity and \eqref{eq:high-graph-tail} bound the
operator norm of $T_{\Lambda,g}$ on the selected space by
$C\Lambda^{\sigma-M(1-\sigma/s)/2}$. Its dimension is
$O(\Lambda)$. Multiply the squared operator bound by this
dimension and by $\|G_\Lambda^{-1}\|$.
\end{proof}

The corollary isolates the arithmetic inputs: a fixed exponent
$\sigma<1/2$, an exact band-tail identity, and a polynomial inverse
Gram bound. It makes no cancellation assumption for ordinary
M\"obius sums. It proves no estimate for the fixed-threshold
$1/2$ selection, which is a different cutoff. That selection, and
every other fixed threshold in $(0,1)$, will be treated in
Theorem~\ref{nh:theorem}.

